% problem_id: S001-theorem-final-bootstrap
% source_id: S001 (Stacks Project)
% source_file: bootstrap.tex
% statement_locator: bootstrap.tex lines 1445--1470
% proof_locator: bootstrap.tex lines 1472--1652
% env: theorem
% status: text-extracted-verbatim (difficulty judgement pending, written by hand)
%
% The statement and proof below are the author's own text, extracted verbatim.
% Nothing is paraphrased, shortened, or supplemented.

\begin{theorem}
\label{theorem-final-bootstrap}
Let $S$ be a scheme.
Let $F : (\Sch/S)_{fppf}^{opp} \to \textit{Sets}$ be a functor.
Any one of the following conditions implies that $F$ is an algebraic space:
\begin{enumerate}
\item $F = U/R$ where $(U, R, s, t, c)$ is a groupoid in algebraic spaces
over $S$ such that $s, t$ are flat and locally of finite presentation, and
$j = (t, s) : R \to U \times_S U$ is an equivalence relation,
\item $F = U/R$ where $(U, R, s, t, c)$ is a groupoid scheme
over $S$ such that $s, t$ are flat and locally of finite presentation, and
$j = (t, s) : R \to U \times_S U$ is an equivalence relation,
\item $F$ is a sheaf and there exists an algebraic space $U$ and a morphism
$U \to F$ which is representable by algebraic spaces,
surjective, flat and locally of finite presentation,
\item $F$ is a sheaf and there exists a scheme $U$ and a morphism
$U \to F$ which is representable by algebraic spaces or schemes,
surjective, flat and locally of finite presentation,
\item $F$ is a sheaf, $\Delta_F$ is representable by algebraic spaces,
and there exists an algebraic space $U$ and a morphism $U \to F$ which is
surjective, flat, and locally of finite presentation, or
\item $F$ is a sheaf, $\Delta_F$ is representable,
and there exists a scheme $U$ and a morphism $U \to F$ which is
surjective, flat, and locally of finite presentation.
\end{enumerate}
\end{theorem}

\begin{proof}
Trivial observations: (6) is a special case of (5) and
(4) is a special case of (3).
We first prove that cases (5) and (3) reduce to case (1).
Namely, by bootstrapping the diagonal
Lemma \ref{lemma-bootstrap-diagonal}
we see that (3) implies (5). In case (5) we set $R = U \times_F U$ which
is an algebraic space by assumption. Moreover, by assumption both
projections $s, t : R \to U$ are surjective, flat and locally of
finite presentation. The map $j : R \to U \times_S U$ is clearly an
equivalence relation. By
Lemma \ref{lemma-surjective-flat-locally-finite-presentation}
the map $U \to F$ is a surjection of sheaves. Thus $F = U/R$
which reduces us to case (1).

\medskip\noindent
Next, we show that (1) reduces to (2).
Namely, let $(U, R, s, t, c)$ be a groupoid in algebraic spaces
over $S$ such that $s, t$ are flat and locally of finite presentation, and
$j = (t, s) : R \to U \times_S U$ is an equivalence relation.
Choose a scheme $U'$ and a surjective \'etale morphism $U' \to U$.
Let $R' = R|_{U'}$ be the restriction of $R$ to $U'$. By
Groupoids in Spaces,
Lemma \ref{spaces-groupoids-lemma-quotient-pre-equivalence-relation-restrict}
we see that $U/R = U'/R'$. Since $s', t' : R' \to U'$ are also
flat and locally of finite presentation (see
More on Groupoids in Spaces,
Lemma \ref{spaces-more-groupoids-lemma-restrict-preserves-type})
this reduces us to the case where $U$ is a scheme.
As $j$ is an equivalence relation we see that $j$ is a monomorphism.
As $s : R \to U$ is locally of finite presentation we see that
$j : R \to U \times_S U$ is locally of finite type, see
Morphisms of Spaces, Lemma \ref{spaces-morphisms-lemma-permanence-finite-type}.
By
Morphisms of Spaces, Lemma
\ref{spaces-morphisms-lemma-monomorphism-loc-finite-type-loc-quasi-finite}
we see that $j$ is locally quasi-finite and separated.
Hence if $U$ is a scheme, then $R$ is a scheme by
Morphisms of Spaces, Proposition
\ref{spaces-morphisms-proposition-locally-quasi-finite-separated-over-scheme}.
Thus we reduce to proving the theorem in case (2).

\medskip\noindent
Assume $F = U/R$ where $(U, R, s, t, c)$ is a groupoid scheme
over $S$ such that $s, t$ are flat and locally of finite presentation, and
$j = (t, s) : R \to U \times_S U$ is an equivalence relation. By
Lemma \ref{lemma-slice-equivalence-relation}
we reduce to that case where $s, t$ are flat,
locally of finite presentation, and locally quasi-finite.
Let $U = \bigcup_{i \in I} U_i$ be an affine open covering
(with index set $I$ of cardinality $\leq$ than the size of $U$ to avoid
set theoretic problems later -- most readers can safely ignore this remark).
Let $(U_i, R_i, s_i, t_i, c_i)$ be the restriction of $R$
to $U_i$. It is clear that $s_i, t_i$ are still flat, locally of finite
presentation, and locally quasi-finite as $R_i$ is the open subscheme
$s^{-1}(U_i) \cap t^{-1}(U_i)$ of $R$
and $s_i, t_i$ are the restrictions of $s, t$ to this open. By
Lemma \ref{lemma-better-finding-opens}
(or the simpler
Spaces, Lemma \ref{spaces-lemma-finding-opens})
the map $U_i/R_i \to U/R$ is representable by open immersions.
Hence if we can show that $F_i = U_i/R_i$ is an algebraic space, then
$\coprod_{i \in I} F_i$ is an algebraic space by
Spaces, Lemma \ref{spaces-lemma-coproduct-algebraic-spaces}.
As $U = \bigcup U_i$ is an open covering it is clear that
$\coprod F_i \to F$ is surjective. Thus
it follows that $U/R$ is an algebraic space, by
Spaces, Lemma \ref{spaces-lemma-glueing-algebraic-spaces}.
In this way we reduce to the case where $U$ is affine and $s, t$ are flat,
locally of finite presentation, and locally quasi-finite and
$j$ is an equivalence.

\medskip\noindent
Assume $(U, R, s, t, c)$ is a groupoid scheme over $S$,
with $U$ affine, such that $s, t$ are flat, locally of finite presentation,
and locally quasi-finite, and $j$ is an equivalence relation.
Choose $u \in U$. We apply
More on Groupoids in Spaces,
Lemma \ref{spaces-more-groupoids-lemma-quasi-splitting-affine-scheme}
to $u \in U, R, s, t, c$. We obtain an affine scheme $U'$, an \'etale
morphism $g : U' \to U$, a point $u' \in U'$ with $\kappa(u) = \kappa(u')$
such that the restriction $R' = R|_{U'}$ is quasi-split over $u'$.
Note that the image $g(U')$ is open as $g$ is \'etale and contains $u$.
Hence, repeatedly applying the lemma, we can find finitely many
points $u_i \in U$, $i = 1, \ldots, n$,
affine schemes $U'_i$, \'etale morphisms $g_i : U_i' \to U$, points
$u'_i \in U'_i$ with $g(u'_i) = u_i$ such that (a) each
restriction $R'_i$ is quasi-split over some point in $U'_i$ and
(b) $U = \bigcup_{i = 1, \ldots, n} g_i(U'_i)$.
Now we rerun the last part of the argument in the preceding paragraph:
Using
Lemma \ref{lemma-better-finding-opens}
(or the simpler
Spaces, Lemma \ref{spaces-lemma-finding-opens})
the map $U'_i/R'_i \to U/R$ is representable by open immersions.
If we can show that $F_i = U'_i/R'_i$ is an algebraic space, then
$\coprod_{i \in I} F_i$ is an algebraic space by
Spaces, Lemma \ref{spaces-lemma-coproduct-algebraic-spaces}.
As $\{g_i : U'_i \to U\}$ is an \'etale covering
it is clear that $\coprod F_i \to F$ is surjective. Thus
it follows that $U/R$ is an algebraic space, by
Spaces, Lemma \ref{spaces-lemma-glueing-algebraic-spaces}.
In this way we reduce to the case where $U$ is affine and $s, t$ are flat,
locally of finite presentation, and locally quasi-finite,
$j$ is an equivalence, and $R$ is quasi-split over $u$ for some
$u \in U$.

\medskip\noindent
Assume $(U, R, s, t, c)$ is a groupoid scheme over $S$,
with $U$ affine, $u \in U$ such that $s, t$ are flat, locally
of finite presentation, and locally quasi-finite and
$j = (t, s) : R \to U \times_S U$ is an equivalence relation
and $R$ is quasi-split over $u$. Let $P \subset R$ be a quasi-splitting
of $R$ over $u$. By
Lemma \ref{lemma-divide-subgroupoid}
we see that $(U, R, s, t, c)$ is the restriction of a groupoid
$(\overline{U}, \overline{R}, \overline{s}, \overline{t}, \overline{c})$
by a surjective finite locally free morphism $U \to \overline{U}$ such that
$P = U \times_{\overline{U}} U$. Note that $s$ admits a factorization
$$
R = U \times_{\overline{U}, \overline{t}} \overline{R}
\times_{\overline{s}, \overline{U}} U
\xrightarrow{\text{pr}_{23}}
\overline{R} \times_{\overline{s}, \overline{U}} U
\xrightarrow{\text{pr}_2} U
$$
The map $\text{pr}_2$ is the base change of $\overline{s}$, and
the map $\text{pr}_{23}$ is a base change of the surjective finite locally
free map $U \to \overline{U}$. Since $s$ is flat, locally
of finite presentation, and locally quasi-finite and since $\text{pr}_{23}$
is surjective finite locally free (as a base change of such), we
conclude that $\text{pr}_2$ is flat, locally
of finite presentation, and locally quasi-finite by
Descent, Lemmas
\ref{descent-lemma-flat-fpqc-local-source} and
\ref{descent-lemma-locally-finite-presentation-fppf-local-source} and
Morphisms, Lemma \ref{morphisms-lemma-quasi-finite-local-source}.
Since $\text{pr}_2$ is the base change of the morphism
$\overline{s}$ by $U \to \overline{U}$ and $\{U \to \overline{U}\}$
is an fppf covering we conclude $\overline{s}$ is
flat, locally of finite presentation, and locally quasi-finite, see
Descent, Lemmas \ref{descent-lemma-descending-property-flat},
\ref{descent-lemma-descending-property-locally-finite-presentation}, and
\ref{descent-lemma-descending-property-quasi-finite}. The same goes
for $\overline{t}$. Consider the commutative diagram
$$
\xymatrix{
U \times_{\overline{U}} U \ar@{=}[r] \ar[rd] & P \ar[r] \ar[d] & R \ar[d] \\
& \overline{U} \ar[r]^{\overline{e}} & \overline{R}
}
$$
It is a general fact about restrictions that the outer four corners
form a cartesian diagram. By the equality we see the inner square is
cartesian. Since $P$ is open in $R$ (by definition of a quasi-splitting)
we conclude that $\overline{e}$ is an open immersion by
Descent, Lemma \ref{descent-lemma-descending-property-open-immersion}.
An application of
Groupoids,
Lemma \ref{groupoids-lemma-quotient-pre-equivalence-relation-restrict}
shows that $U/R = \overline{U}/\overline{R}$. Hence we have reduced to
the case where $(U, R, s, t, c)$ is a groupoid scheme over $S$,
with $U$ affine, $u \in U$ such that $s, t$ are flat, locally
of finite presentation, and locally quasi-finite and
$j = (t, s) : R \to U \times_S U$ is an equivalence relation
and $e : U \to R$ is an open immersion!

\medskip\noindent
But of course, if $e$ is an open immersion and
$s, t$ are flat and locally of finite presentation
then the morphisms $t, s$ are \'etale.
For example you can see this by applying
More on Groupoids, Lemma \ref{more-groupoids-lemma-sheaf-differentials}
which shows that $\Omega_{R/U} = 0$ which in turn implies
that $s, t : R \to U$ is G-unramified (see
Morphisms, Lemma \ref{morphisms-lemma-unramified-omega-zero}),
which in turn implies that $s, t$ are \'etale (see
Morphisms, Lemma \ref{morphisms-lemma-flat-unramified-etale}).
And if $s, t$ are \'etale then finally $U/R$ is an algebraic
space by
Spaces, Theorem \ref{spaces-theorem-presentation}.
\end{proof}
